% !TEX root = ../main.tex
\section{Resultados}
\label{sec:resultados}

\subsection{Efeito do controle de qualidade dos dados}
\label{sec:preproc}

Antes de comparar arquiteturas entre si, é preciso isolar um efeito que domina os
resultados: o pré-processamento. Empilhar tratamentos uns sobre os outros (como em
versões anteriores deste texto) não permite separar o que cada etapa
individualmente contribui --- um ganho visto na combinação ``solo + filtro +
outliers'' pode vir inteiro de um único desses três, com os outros dois sendo
irrelevantes ou até prejudiciais. Para resolver isso, cada tratamento foi testado
\textbf{isoladamente}, sempre contra a mesma nuvem crua (\texttt{raw}), sem
empilhar com os demais: projeção do solo sozinha, filtro $\leq$1,3~m sozinho
(sobre o deslocamento mínimo de Z, sem o relevo local), remoção de outliers
geométricos sozinha, e correção de gap temporal sozinha no alvo. Uma sexta
coluna, ``tudo combinado'', aplica os quatro tratamentos juntos --- o pipeline
completo. A Tabela~\ref{tab:comparacao} reporta as seis colunas para os sete
modelos.

\textbf{Remoção de outliers geométricos é, disparado, o tratamento que mais
importa} --- e o único que, sozinho, já entrega quase todo o ganho visto na
combinação completa. Nos três modelos com as duas variantes disponíveis para
comparação direta (ABA, Random Forest e CNN 2D), o outlier isolado iguala ou
\textbf{supera} a combinação ``tudo junto'': CNN 2D vai de $R^2 \approx$ 0,25
(nuvem crua) para \textbf{0,55} só removendo outliers --- mais alto que os 0,48 da
combinação completa (solo+filtro+outliers+gap). Random Forest (0,53 isolado vs.
0,51 combinado) e ABA (0,51 vs. 0,49) têm o mesmo padrão. Isso indica que solo,
filtro e gap, no pipeline combinado, não estavam contribuindo --- ou estavam
discretamente atrapalhando --- e o ganho inteiro vinha da geometria das parcelas.
O GBR de 32 quantis confirma a mesma magnitude isolada ($R^2$ = 0,51 só com
outliers, Tabela~\ref{tab:comparacao}), ainda que sem uma variante ``tudo
combinado'' treinada para comparação direta.

\textbf{Projeção do solo e filtro $\leq$1,3~m, isolados, têm efeito pequeno e
inconsistente entre arquiteturas.} Só a projeção do solo (sem filtro) move o $R^2$
em poucos centésimos na maioria dos modelos (ex.: ABA 0,320$\rightarrow$0,328,
Random Forest 0,283$\rightarrow$0,286) --- exceção é a CNN 3D (voxel), que sai de
$R^2$ negativo ($-$0,13) para positivo (0,09), sugerindo que a estrutura vertical
bruta (Z absoluto, sem corrigir relevo) é especialmente ruim pra arquiteturas que
processam o volume 3D diretamente. O filtro $\leq$1,3~m isolado (sobre o mesmo
deslocamento mínimo de Z do baseline cru, sem acionar a projeção de relevo
completa) tem efeito ainda menor: piora levemente ABA, Random Forest, o modelo de
alturas ordenadas e a PointNet ($R^2$ cai alguns centésimos em cada caso ---
PointNet 0,125$\rightarrow$0,118) e ajuda discretamente só a CNN 2D e a CNN 3D,
que parecem se beneficiar de qualquer redução no ruído vertical bruto, mesmo
pequena.

\textbf{Correção de gap temporal isolada é, na maioria dos modelos, mais fraca que
não fazer nada.} Em cinco dos seis modelos com baseline cru comparável, o gap
isolado fica abaixo da nuvem crua sem nenhum tratamento --- exceção é a CNN 3D,
onde o gap isolado ($R^2 = -0,05$) é menos ruim que a própria nuvem crua
($-0,13$), embora ainda negativo. Isso contrasta com o papel do gap na combinação
completa, onde ele aparece ``ajudando'' --- mas como visto acima, o ganho ali vem
majoritariamente da remoção de outliers que o acompanha, não do gap em si. A
correção de gap é mantida no pipeline por ser metodologicamente correta (alinhar
DBH/altura do inventário à data do voo LiDAR, evitando viés sistemático de
subestimar biomassa quando o LiDAR é mais recente que o censo), não porque eleve o
$R^2$.

Em suma: \textbf{a qualidade da curadoria geométrica das parcelas (outliers)
domina a escolha de arquitetura em magnitude de efeito} --- nenhuma combinação de
solo, filtro ou gap chega perto do ganho de simplesmente descartar as parcelas de
geometria ruim. Este é o primeiro resultado do trabalho, e não é sutil.

\begin{table}[H]
  \centering
  \caption{Comparação entre modelos e variantes de pré-processamento (RMSE em
  Mg/ha; rRMSE em \%; $R^2$ adimensional; maior $R^2$ e menor RMSE/rRMSE são
  melhores). Colunas ``Só X'' são tratamentos isolados contra a mesma nuvem crua;
  ``Tudo combinado'' empilha os quatro juntos.}
  \label{tab:comparacao}
  % PLACEHOLDER — os valores por célula são gerados pelos scripts de treino
  % (data/processed/06_model/cv_results_*.csv) e exibidos no dashboard.
  % Exporte a tabela final e cole aqui no formato booktabs, ex.:
  % \begin{tabular}{lrrrrrr}
  %   \toprule
  %   Modelo & Nuvem crua & Só solo & Só filtro & Só outliers & Só gap & Tudo combinado \\
  %   \midrule
  %   ... & ... \\
  %   \bottomrule
  % \end{tabular}
  \figplaceholder{[PLACEHOLDER] Tabela 1 --- preencher com os valores de
  RMSE/$R^2$/rRMSE por modelo (7 linhas) $\times$ variante (6 colunas), exportados
  de \texttt{data/processed/06\_model/cv\_results\_*.csv}.}
  \fonte{elaborado pelo autor (2026), a partir das validações cruzadas do
  pipeline de treino}
\end{table}

\subsection{Comparação entre arquiteturas, na melhor configuração}
\label{sec:arch}

Restringindo a comparação à melhor variante de CADA modelo individualmente
($n = 325$) --- e não assumindo de antemão qual tratamento essa variante usa, já
que a Seção~\ref{sec:preproc} mostrou que cada tratamento isolado tem efeito muito
diferente ---, a variante líder é \textbf{``só outliers'' (isolada, sem
solo/filtro/gap) para cinco dos seis modelos}; CNN 3D e PointNet são as exceções,
precisando do pipeline combinado completo para atingir seu melhor resultado. A
ordem que emerge é: \textbf{CNN 2D sobre rasters de altura ($R^2$ = 0,547) $>$
Random Forest (0,532) $>$ GBR sobre a distribuição completa de 1024 alturas
(0,526) $>$ ABA (0,514) $\approx$ GBR sobre 32 quantis de altura (0,510) $>$
CNN 3D sobre voxels (0,455, pipeline combinado) $\gg$ PointNet sobre a nuvem crua
(0,212, pipeline combinado, único caso em que a variante gap-corrigida vence).}
(Esta ordem usa a fórmula alométrica de referência M1; a
Seção~\ref{sec:allom-sens} mostra que ela não é robusta à troca para M3 --- o
Random Forest assume a liderança, embora essa análise de sensibilidade ainda não
tenha sido refeita sobre as variantes isoladas identificadas aqui, ver ressalva ao
final da Seção~\ref{sec:allom-sens}.)

Essa ordem corrige uma leitura anterior deste texto, que agrupava Random Forest e
CNN 3D como equivalentes: usando a variante correta por modelo, a CNN 3D fica
\textbf{abaixo} de todos os modelos estruturais/estatísticos (ABA, GBR em ambas as
representações), não empatada com o Random Forest. A CNN 2D ainda vence por uma
margem modesta sobre Random Forest (diferença de $R^2 \approx$ 0,015 --- dentro da
variabilidade esperada de um k-fold de 5 dobras; a Seção~\ref{sec:confounders}
reporta um teste pareado formal entre CNN 2D, Random Forest e CNN 3D, cada um na
sua própria melhor variante, que não rejeita a hipótese de diferença nula entre as
três), o que sugere que a estrutura \textbf{horizontal} do dossel capturada pelo
raster (clareiras, agrupamento de copas --- informação que o ABA e o GBR de
alturas descartam por trabalhar só com a distribuição vertical agregada) carrega
sinal real, mas marginal frente à estrutura vertical já capturada pelas abordagens
mais simples.

O resultado mais informativo da tabela, porém, é o \textbf{fracasso relativo da
PointNet}: mesmo na sua melhor configuração, $R^2 \approx$ 0,21 --- pior que
\textit{todos} os outros cinco modelos, incluindo o GBR mais simples da
comparação. A PointNet recebe a representação mais rica em princípio (a nuvem de
pontos crua, sem nenhuma agregação), mas é também a que mais depende de volume de
dados para aprender uma função de agregação útil via max-pooling --- e 325
parcelas de treino são, evidentemente, poucas para isso.

\begin{figure}[H]
  \centering
  % PLACEHOLDER — substituir pelo scatter predito × observado:
  % \includegraphics[width=0.7\textwidth]{figura2_pred_obs.png}
  \figplaceholder{[PLACEHOLDER] Figura 2 --- Predito $\times$ observado, melhor
  modelo (CNN 2D, só outliers, isolado).}
  \caption{Predito $\times$ observado, melhor modelo (CNN 2D, só outliers,
  isolado). A dispersão em torno da diagonal 1:1 mostra o padrão típico de um
  modelo com $R^2$ moderado: tendência central capturada, mas variância residual
  substancial, sobretudo nas parcelas de AGB mais alta --- o regime onde erros
  absolutos (Mg/ha) pesam mais para qualquer aplicação de contabilidade de
  carbono.}
  \label{fig:pred-obs}
  \fonte{elaborado pelo autor (2026)}
\end{figure}

\subsection{Ablação de dimensionalidade: o achado central}
\label{sec:ablation}

A comparação mais direta com a pergunta de pesquisa deste trabalho não é entre
arquiteturas, mas entre \textbf{quantidade de informação de altura} fornecida ao
mesmo algoritmo (Gradient Boosting). A Figura~\ref{fig:ablation} mostra o $R^2$ de
validação cruzada em função do número de quantis de altura usados como atributo,
de 8 a 256, com o vetor completo de 1024 alturas ordenadas como referência.

O padrão contraria a intuição de que ``mais informação bruta é sempre melhor'':
\textbf{o vetor completo de 1024 alturas ($R^2 \approx$ 0,39--0,42) é superado já
pela representação mais grosseira da ablação, 8 quantis ($R^2 \approx$ 0,43)} ---
menos de 1\% da dimensionalidade original captura mais sinal preditivo, não menos.
O ganho de 8 para 16 quantis é real e consistente ($R^2$ 0,426 $\rightarrow$
0,442); daí em diante, até 64 quantis, o desempenho \textbf{satura} ($R^2$ oscila
entre 0,441 e 0,444, dentro do ruído do k-fold); só volta a subir de forma
perceptível em 256 quantis ($R^2 \approx$ 0,455) --- um ganho pequeno pelo custo de
16$\times$ mais atributos que o ponto de saturação.

\begin{figure}[H]
  \centering
  % PLACEHOLDER — substituir pela curva de ablação:
  % \includegraphics[width=0.7\textwidth]{figura3_ablation.png}
  \figplaceholder{[PLACEHOLDER] Figura 3 --- $R^2$ de validação cruzada vs. número
  de quantis de altura (escala log em x).}
  \caption{$R^2$ de validação cruzada vs. número de quantis de altura (escala log
  em $x$).}
  \label{fig:ablation}
  \fonte{elaborado pelo autor (2026)}
\end{figure}

A leitura estatística é direta: com 325 parcelas de treino e um vetor de 1024
atributos altamente correlacionados entre si (alturas vizinhas de uma distribuição
ordenada), o modelo tem mais graus de liberdade que sinal identificável --- o
excesso de dimensão é absorvido como ruído de treino (sobreajuste), não como
capacidade preditiva adicional. Reduzir para algumas dezenas de quantis não é
perda de informação relevante; é remoção de ruído correlacionado que atrapalhava a
otimização.

\subsection{O que o modelo interpretável revela}
\label{sec:interpret}

A única abordagem desta comparação que expõe explicitamente \textbf{quais
propriedades da floresta} sustentam a previsão é o ABA (GBR sobre 37 métricas
estruturais nomeadas). A Figura~\ref{fig:importance} mostra a importância das 10
variáveis mais relevantes no modelo ABA na sua variante de melhor desempenho ---
outliers isolado (Seção~\ref{sec:preproc}/\ref{sec:arch}), não o pipeline
combinado usado em versões anteriores deste texto.

\begin{figure}[H]
  \centering
  % PLACEHOLDER — substituir pelo gráfico de importância de variáveis:
  % \includegraphics[width=0.7\textwidth]{figura4_importance.png}
  \figplaceholder{[PLACEHOLDER] Figura 4 --- Importância de variáveis, modelo ABA
  (top 10).}
  \caption{Importância de variáveis, modelo ABA (top 10).}
  \label{fig:importance}
  \fonte{elaborado pelo autor (2026)}
\end{figure}

A \textbf{entropia da densidade vertical} (\texttt{vfd\_entropy}) domina
isoladamente a importância agregada ($\approx$0,34 do total), mais que o dobro da
segunda colocada --- um resultado que muda a leitura anterior deste texto,
centrada nos percentis de altura. Essa métrica mede quão uniformemente distribuída
está a folhagem ao longo do perfil vertical (entropia de Shannon sobre a densidade
de pontos por faixa de altura): um dossel com muitos estratos de folhagem bem
distribuídos tem entropia alta; um dossel concentrado numa única camada (ex.:
topo denso, meio vazio) tem entropia baixa. Apenas o percentil p75 ($\approx$0,21)
mantém posição de destaque entre os percentis de altura --- os demais (p50, p25,
p10, p05, que dominavam na variante combinada) caem para a metade inferior da
tabela. As demais variáveis relevantes --- densidade de pontos (\texttt{pt\_density}),
rugosidade convolucional do dossel (\texttt{lap\_roughness}, derivada de um
Laplaciano 3$\times$3 sobre o CHM), dispersão de altura (\texttt{h\_std}) e fração
de clareiras (\texttt{gap\_fraction}) --- carregam informação complementar sobre
heterogeneidade estrutural (vertical e horizontal) que a altura típica sozinha não
captura, e ajudam a explicar por que o raster 2D (que também enxerga
heterogeneidade horizontal diretamente, sem precisar de uma métrica resumo como
\texttt{vfd\_entropy}) supera marginalmente o ABA puro.

\subsection{Posição frente à literatura e ao GEDI}
\label{sec:gedi}

O melhor resultado desta tese ($R^2$ = 0,547; rRMSE $\approx$ 41,7\%; CNN 2D,
parcela, variante só outliers isolada) situa-se \textbf{abaixo} dos trabalhos de
referência em LiDAR aéreo tropical com metodologia mais madura --- Asner \&
Mascaro (2014), $R^2$ = 0,90 sobre 804 parcelas de 1~ha em quatro países tropicais
e Havaí; Longo et al. (2016), $R^2$ = 0,70 sobre 359 parcelas amazônicas --- e
também \textbf{abaixo} de Bispo et al. (2020): $R^2$ = 0,93 / rRMSE = 13\% para o
modelo só-LiDAR ($R^2$ = 0,89 / rRMSE = 43\% no modelo multissensor). \textit{Nota:
Bispo et al. (2020) estuda o Cerrado brasileiro (bacia do Rio Vermelho), não a
Amazônia --- versões anteriores deste texto atribuíam o estudo à Amazônia e
citavam $R^2$ = 0,58 / rRMSE = 36,0\%, que não correspondem à fonte; ambos os
erros foram corrigidos após verificação adversarial direta contra o artigo
original.} Frente ao produto orbital \textbf{GEDI L4A} (\textit{footprint}
$\sim$25~m), a validação independente em savanas do sul da África (Li et al., 2024
--- não Naidoo et al., que é coautor, não autor líder; ver correção na seção de
Referências) reporta $R^2 \approx$ 0,42 --- nesse critério específico, o resultado
desta tese é \textbf{competitivo}, apesar de operar em escala de parcela (mais
exigente) em vez de \textit{footprint} de 25~m. d'Oliveira et al. (2012), com
suporte espacial menor (50 parcelas de 0,25~ha, Acre) e instrumentação mais
antiga, reporta $R^2$ = 0,63--0,72 --- o mais baixo do conjunto de referências ---
e foi mantido na lista de referências sem entrar nesta comparação direta por essas
duas diferenças de desenho.

Essas comparações exigem duas ressalvas, detalhadas na Seção~\ref{sec:lit-check}:
(i) as escalas espaciais não são diretamente comparáveis (parcela vs.
\textit{footprint} vs. grade de 1~km); (ii), mais importante, os números desta
tese vêm de k-fold aleatório, enquanto boa parte da literatura citada usa
validação cruzada por sítio ou validação independente --- um padrão de
generalização mais exigente que o empregado aqui (Seção~\ref{sec:spatial-validation}).
A verificação adversarial dos números da literatura --- conferência linha a linha
contra a fonte primária, antes pendente --- foi concluída; ver
Seção~\ref{sec:lit-check} para o detalhe por fonte.

Quanto ao LiDAR orbital da NASA --- GEDI e ICESat-2 ---, não foi encontrada, mesmo
após busca e verificação direta contra a fonte primária, nenhuma validação de
$R^2$/RMSE calculada isoladamente para a Amazônia brasileira (nem do GEDI, nem do
ICESat-2) com desenho comparável a parcela ou k-fold aleatório. O benchmark
publicado mais próximo é Duncanson et al. (2022): o modelo oficial de biomassa do
GEDI, no estrato \textit{Evergreen Broadleaf Trees} --- América do Sul (classe
dominante na Amazônia, mas estrato continental, não um recorte amazônico), reporta
$R^2$ = 0,66 e \%RMSE = 42,2\% por validação cruzada geográfica --- portanto nem
k-fold aleatório, nem especificamente amazônico, mas é o melhor proxy disponível
para ``GEDI num bioma análogo à Amazônia'', na mesma faixa de rRMSE desta tese
(41,7\%). Já o ICESat-2 (ATL08), quando validado contra LiDAR aerotransportado em
floresta tropical (Mesoamérica, Gabão e Guiana Francesa --- nunca a bacia
amazônica isolada) (Fernandez-Diaz; Velikova; Glennie, 2022), degrada fortemente frente ao desempenho em floresta boreal:
erro de altura de dossel de 5--7~m em floresta tropical contra 3,19~m em condições
boreais ideais (Liu, Gong \& Shi, 2026) --- reforçando que mesmo o LiDAR orbital
mais recente da NASA perde precisão em dossel tropical denso, um argumento a favor
da abordagem aerotransportada adotada nesta tese. O trabalho de deep learning mais
próximo publicado em 2023--2025 usando a mesma linhagem de dados (LiDAR
aerotransportado do programa NASA CMS Brasil/EBA) na Amazônia é Wagner et al.
(2025, pré-print): uma U-Net treinada com esse mesmo tipo de levantamento para
mapear altura de dossel (não AGB) a partir de imagens Planet NICFI, com
MAE = 3,68~m em validação espacialmente estratificada --- uma comparação
metodológica (arquitetura de rede) mais honesta que uma comparação direta de
métrica de erro, já que a variável-alvo difere da desta tese.
